\section{Problem \& Opportunity}
\label{sec:problem_opportunity}

\subsection{Renewable Generation Variability}
\label{subsec:variability}

The global transition to clean power systems relies heavily on solar and wind generation. However, unlike dispatchable thermal or hydroelectric generation, renewable output is intrinsically variable and governed by transient meteorological conditions. Solar power fluctuates with solar irradiance, cloud transients, ambient temperature, and diurnal cycles. Wind generation varies with atmospheric pressure fronts and turbulent wind speeds, scaling non-linearly with the cube of velocity.

This inherent stochasticity introduces severe operational uncertainty: \emph{power grid operators must schedule generation and balance supply and demand in advance, based on imperfect information about weather conditions that will unfold 24 to 72 hours later.}

\subsection{The Imbalance Challenge: Over- and Under-Generation}
\label{subsec:imbalance_challenge}

When renewable generation deviates from expected schedules, two critical grid conditions emerge:
\begin{itemize}[leftmargin=1.5em, itemsep=0.2em]
    \item \textbf{Under-Generation (Supply Shortfall):} Renewable output falls below scheduled baselines or net demand. Without advance warning, grid operators must rapidly dispatch expensive, carbon-intensive fossil peaker plants or risk under-frequency events and emergency load shedding.
    \item \textbf{Over-Generation (Energy Surplus):} Renewable generation exceeds immediate local demand and available transmission hosting capacity. If storage or export channels cannot absorb the excess, operators are forced to curtail clean generation, resulting in the permanent loss of zero-marginal-cost renewable energy.
\end{itemize}

\subsection{Why Forecasting Alone Is Insufficient}
\label{subsec:forecast_gap}

The HackOut'26 challenge mandates predicting 24--72 hour solar and wind generation using weather, historical generation, and site parameters, while identifying surplus/shortfall periods and recommending grid actions (curtailment, storage dispatch, backup activation).

Conventional utility platforms address only the first part of this mandate by providing a deterministic point prediction (e.g., \emph{``Tomorrow at 14:00: Expected Output = 80~MW''}). This solitary number creates a critical \textbf{forecast-to-decision gap}:
\begin{itemize}[leftmargin=1.5em, itemsep=0.15em]
    \item It fails to communicate \textbf{uncertainty} (how wide is the statistical confidence interval?).
    \item It omits \textbf{risk quantification} (what is the probability of a supply deficit during the evening ramp?).
    \item It provides no \textbf{operational guidance} (should battery storage charge now or reserve capacity for a subsequent surplus?).
\end{itemize}

\subsection{Target Stakeholders and Core Opportunity}
\label{subsec:target_stakeholders}

Grid operators, utility companies, renewable plant owners, and energy traders all require forward operational visibility. The primary opportunity of GridSense AI is to convert passive generation forecasts into an active, uncertainty-aware **decision-support system** that bridges meteorological predictions and control-room actions.
